% Λύση του 42ου προβλήματος της ενότητας «Άλυτα Προβλήματα».
\subsection*{Λύση Προβλήματος 42 --- Ορίζοντας ραδιοζεύξης}

\begin{alist}
  \item Η οριακή ακτίνα εκπομπής αγγίζει την επιφάνεια της Γης στο
        σημείο $T$. Επομένως είναι εφαπτομένη στον κύκλο που
        προσεγγίζει τη Γη.

        Η ακτίνα στο σημείο επαφής είναι κάθετη στην εφαπτομένη. Άρα
        \[
          OT\perp AT.
        \]
        Το τρίγωνο $OAT$ είναι ορθογώνιο στο $T$.

        Έχουμε
        \[
          OT=R,
          \qquad
          OA=R+h,
          \qquad
          AT=d.
        \]
        Από το Πυθαγόρειο Θεώρημα:
        \[
          d^2+R^2=(R+h)^2.
        \]
        Άρα
        \[
          {d=\sqrt{(R+h)^2-R^2}}.
        \]

  \item Αντικαθιστούμε
        \[
          R=6378\si{km},
          \qquad
          h=0{,}1\si{km}.
        \]
        Τότε
        \[
          d=\sqrt{(6378+0{,}1)^2-6378^2}.
        \]
        Υπολογίζουμε:
        \[
          d=\sqrt{2\cdot6378\cdot0{,}1+0{,}1^2}.
        \]
        Επομένως
        \[
          d=\sqrt{1275{,}6+0{,}01}
          =
          \sqrt{1275{,}61}
          \approx35{,}7.
        \]
        Άρα η ευθύγραμμη απόσταση μέχρι τον ορίζοντα είναι περίπου
        \[
          {35{,}7\si{km}}.
        \]

  \item Ένας δέκτης που βρίσκεται σε ευθύγραμμη απόσταση
        $40\si{km}$ από την κορυφή της κεραίας είναι πιο μακριά από την
        απόσταση του ορίζοντα, αφού
        \[
          40>35{,}7.
        \]
        Άρα, στο μοντέλο χωρίς ατμοσφαιρική διάθλαση, η καμπυλότητα της
        Γης εμποδίζει την απευθείας οπτική επαφή.

        Επομένως ο δέκτης
        \[
          {\text{δεν μπορεί να έχει απευθείας οπτική επαφή}}
        \]
        με την κεραία.
\end{alist}
